% Índice breve de secciones numeradas para las aperturas de capítulo.
% Se obtiene del .toc existente y se actualiza al compilar con latexmk.
\newlength{\indiceCapituloAncho}

\makeatletter
\def\miniindice@clave#1.#2\@nil{\def\miniindice@capitulo{#1}}
\long\def\miniindice@guardar#1#2#3\@nil{%
  \miniindice@clave#2.\@nil
  \ifcsname miniindice@lista@\miniindice@capitulo\endcsname\else
    \expandafter\gdef\csname miniindice@lista@\miniindice@capitulo\endcsname{}%
    \expandafter\gdef\csname miniindice@total@\miniindice@capitulo\endcsname{0}%
  \fi
  \expandafter\xdef\csname miniindice@total@\miniindice@capitulo\endcsname{%
    \number\numexpr\csname miniindice@total@\miniindice@capitulo\endcsname+1\relax}%
  \expandafter\g@addto@macro\csname miniindice@lista@\miniindice@capitulo\endcsname
    {\filaIndiceCapitulo{#2}{#3}{#1}}%
}
\long\def\miniindice@leer#1#2#3\@nil{%
  \ifx#2\numberline
    \miniindice@guardar{#1}#3\@nil
  \fi
}
\newcommand{\cargarIndicesCapitulos}{%
  \begingroup
    \makeatletter
    \let\babel@toc\@gobbletwo
    \long\def\contentsline##1##2##3##4{%
      \ifstrequal{##1}{section}{\miniindice@leer{##4}##2\@nil}{}%
    }%
    \InputIfFileExists{\jobname.toc}{}{}%
  \endgroup
}
\AtBeginDocument{\cargarIndicesCapitulos}

\newcount\miniindice@posicion
\newcount\miniindice@desde
\newcount\miniindice@hasta
\newcount\miniindice@mitad
\newcommand{\filaIndiceCapitulo}[3]{%
  \advance\miniindice@posicion\@ne
  \ifnum\miniindice@posicion>\miniindice@desde
    \ifnum\miniindice@posicion>\miniindice@hasta\else
      \noindent\hangindent=2.5em\hangafter=1
      \makebox[2.5em][l]{\hyperlink{#3}{\textcolor{rojoCIMAT}{\bfseries#1}}}%
      \hyperlink{#3}{#2}\par\vspace{4.5pt}%
    \fi
  \fi
}
\newcommand{\miniindice@columna}[2]{%
  \begin{minipage}[t]{\dimexpr(\indiceCapituloAncho-24pt)/2\relax}
    \raggedright\parindent=0pt\parskip=0pt
    \miniindice@posicion=0
    \miniindice@desde=#1\relax
    \miniindice@hasta=#2\relax
    \csname miniindice@lista@\thechapter\endcsname
    \vspace{-4.5pt}%
  \end{minipage}%
}
\newcommand{\indiceCapitulo}{%
  \begingroup
    \normalfont\fontsize{10}{12.8}\selectfont
    \color{grisCIMAT!75!black}%
    \begin{minipage}{\indiceCapituloAncho}
      \ifcsname miniindice@lista@\thechapter\endcsname
        \noindent\textcolor{grisCIMAT!40!white}{\rule{\linewidth}{.35pt}}\par
        \vspace{9pt}%
        \miniindice@mitad=\csname miniindice@total@\thechapter\endcsname\relax
        \advance\miniindice@mitad\@ne\divide\miniindice@mitad by 2
        \noindent\miniindice@columna{0}{\miniindice@mitad}\hfill
        \miniindice@columna{\miniindice@mitad}{\csname miniindice@total@\thechapter\endcsname}%
      \fi
    \end{minipage}%
  \endgroup
}
\makeatother
